\section{引言}

理解诸如质子这样的强子的内部三维结构，是核物理与粒子物理的一项重要目标。在这方面，横动量依赖（transverse momentum-dependent, TMD）的部分子分布函数（TMDPDFs）~\cite{Collins:1981uk,Collins:1981va}扮演着重要角色，因为它们刻画了强子内在的横向
部分子结构。在 QCD 因子化定理的框架下，这些分布也是描述多尺度与非遍举过程的关键要素，例如电弱规范玻色子或希格斯玻色子的 Drell-Yan 产生，以及小横动量的半遍举深度非弹性散射。因此，过去几十年它们受到了广泛关注（综述见文献~\cite{Angeles-Martinez:2015sea}）。未来几十年，JLab 12 GeV~\cite{Dudek:2012vr}、美国~\cite{Accardi:2012qut,AbdulKhalek:2021gbh}与中国~\cite{Anderle:2021wcy}的电子离子对撞机预期将给出更精确的实验测量。


与编码强子内部分子动量概率密度的 TMDPDFs 不同，横动量依赖波函数（TMDWFs）为强子的部分子结构提供了一种概率幅描述，由此原则上可以计算各种夸克/胶子分布。在涉及横动量的 QCD 因子化中，它们是预言遍举过程物理可观测量——例如对抽取 CKM 矩阵元很有价值并可探寻标准模型之外新物理的重味 $B$ 介子弱衰变~\cite{Keum:2000wi,Lu:2000em}——的最重要成分。然而，由于缺乏对 TMDWFs 的认识，大多数 $B$ 衰变分析使用一维光锥分布振幅（LCDAs）来替代 \cite{Keum:2000wi,Lu:2000em,Nagashima:2002ia}，导致不可控的误差。$B$ 衰变实验测量前所未有的精度~\cite{Cerri:2018ypt}迫切要求获得关于 TMDWFs 的可靠理论知识。

TMDPDFs 与 TMDWFs 的一个共同特征是：它们既依赖于纵向动量分数 $x$，也依赖于部分子间的横向空间间隔。近年来，大量理论工作致力于通过拟合相关实验数据来确定这些量~\cite{Landry:1999an,Landry:2002ix,DAlesio:2014mrz,Sun:2014dqm,Konychev:2005iy,Bacchetta:2017gcc,Scimemi:2017etj,Scimemi:2019cmh,Bacchetta:2019sam}，然而这受到 TMDPDFs 与 TMDWFs 非微扰行为知识不精确的限制。因此，非常希望发展一种从格点 QCD 这类第一性原理方法出发计算它们的途径。

这一目标已在大动量有效理论（large momentum effective theory）~\cite{Ji:2013dva,Ji:2014gla}框架中实现：该理论通过在格点上模拟与时间无关的欧几里得关联函数，为计算光锥关联提供了一条系统化的途径。利用 LaMET 计算各种部分子物理量已取得重大进展；近期综述可参见文献~\cite{Cichy:2018mum,Ji:2020ect}。

LaMET 发展中一项非常重要的成果是：TMDPDFs 与 TMDWFs 可以通过欧几里得的准 TMDPDFs 和准 TMDWFs 以及一个普适软函数（因子）来计算 ~\cite{Ebert:2019okf,Ji:2019sxk,Ji:2019ewn,Ji:2021znw}。文献~\cite{Ji:2019sxk}指出，一类可在格点上计算的、双定域四夸克算符的形状因子，在大动量转移下可通过 QCD 因子化分解为准 TMDWFs、一个普适软（函数）因子与匹配 kernel 的乘积，从而首次实现了在格点上计算普适软函数。于是，光锥 TMD 部分子分布与波函数可以由格点上的四夸克形状因子以及准 TMDPDFs、准 TMDWFs 的数值计算得到~\cite{Ji:2019sxk,Ji:2019ewn}。另一方面，也可以利用 QCD 因子化从准 TMDPDFs 与准 TMDWFs 得到 CS 核。基于这些方案的 CS 核首批结果近来已经发表~\cite{Shanahan:2020zxr,LatticeParton:2020uhz,Schlemmer:2021aij,Li:2021wvl,Shanahan:2021tst}。准 TMDWFs 途径
只需要两点函数的计算，
并有望以相对较小的强子动量逼近光锥极限。

本文基于一圈匹配精度的准 TMDWFs 格点 QCD 分析——在 staggered 海夸克上取 $N_f=2+1+1$ 价 clover 费米子——给出了一项最先进的 CS 核计算。我们采用单一系综：晶格间距 $a\simeq0.12$~fm，体积 $n_s^3\times n_t=48^3\times64$，物理海夸克质量。为了提高信噪比，我们调节轻价夸克质量使得 $m_\pi = 670$ MeV。随后，我们在不同动量 $P^z=2\pi/n_s\times\{8,10,12\}=\{1.72,~2.15,~2.58\}$~GeV 下，通过准 TMDWFs 之比与微扰匹配 kernel 提取 CS 核。
这分别对应洛伦兹增速因子 $\gamma=\{2.57,~3.21,~3.85\}$。本分析纳入了一圈微扰贡献，并分析了来自算符混合、标度依赖的高阶修正以及按 $1/P^z$ 展开的更高幂次修正的系统不确定度，从而改进了先前的工作~\cite{LatticeParton:2020uhz,Li:2021wvl}。

本文其余部分安排如下：
第~\ref{sec:framework}节给出由准 TMDWFs 提取 CS 核的理论框架。准 TMDWFs 与 CS 核的数值结果在第~\ref{sec:numerics}节给出。本文的简要总结见第~\ref{sec:summary}节。分析的更多细节收集于附录中。
